\documentclass[10pt,a4paper]{amsart}
\usepackage[margin=19mm]{geometry}
\usepackage{amsmath,amssymb,mathtools}
\usepackage[scr=rsfs,cal=euler]{mathalfa}
\usepackage{xfrac}
\usepackage[unicode,hidelinks,bookmarks=false]{hyperref}
\newcommand{\ie}{i.e.\ }
\newcommand{\cf}{cf.\ }
\newcommand{\resp}{respectively\ }
\newcommand{\Subsection}{\subsection}
\providecommand{\nobreakdash}{\nobreak\hbox{-}}
\setlength{\emergencystretch}{2em}
\multlinegap0pt
\usepackage{float}
\usepackage{amssymb}
\newtheorem{prop}{Proposition}
\newtheorem{theorem}{Theorem}
\newcommand{\abs}[1]{\left\vert#1\right\vert}
\title{Harmonic K-quasiconformal Koebe functions: construction and application to Pavlovic's problem}
\author{Zhi-Gang Wang, Xiao-Yuan Wang, Antti Rasila, Jia-Le Qiu}
\date{Source: arXiv:2405.19852v5 (CC BY 4.0)}
\begin{document}
\maketitle

\noindent\emph{Source and licence.} Harmonic K-quasiconformal Koebe functions: construction and application to Pavlovic's problem. Version arXiv:2405.19852v5 (CC BY 4.0), distributed by arXiv under the Creative Commons Attribution 4.0 International licence, http://creativecommons.org/licenses/by/4.0/, as recorded in the arXiv licence metadata for this exact version and shown on the abstract page https://arxiv.org/abs/2405.19852v5. Full licence notice: \texttt{licenses/arxiv-2405.19852-CC-BY-4.0.txt}.\par\smallskip

\noindent\textit{Editorial scope.} Counted unit: the article's theorem t1 (source0.tex lines 442--459). The article's complete original proof of it is delivered with it (source0.tex lines 878--906, one \texttt{proof} environment of 29 lines, which the article places away from the statement). No mathematics is written, repaired or re-derived: the statement and the proof are the article's own lines, in the article's own notation, and no retained line is substituted or re-flowed.  Retained uncounted as the source-local context the proof needs: lines 413--421; lines 848--854.  Nothing was omitted from any retained span, so the package carries no omission record.  No retained line contains a citation, so the wrapper carries no bibliography and no bibliography entry.  No retained line contains a figure environment, an image-inclusion command or drawing code, so no image is delivered and the package declares no asset dependency.  Editorial formatting only, in addition: an AMS \texttt{amsart} preamble that restates the article's own declarations and macros used by the retained lines, namely the environments \texttt{prop}, \texttt{theorem} on separate counters, together with the standard packages those lines need.  The retained environments print in this document as Proposition 1, Theorem 1 in order of appearance, whereas the article prints the same items with the numbering of its own full document.  The complete native source of the article as distributed in the arXiv e-print for this exact version is bundled unchanged as \texttt{sources/160002/source0.tex}.
\begin{prop}\label{prop1}
\textit{Suppose that $f=h+\overline{g}\in\widehat{\mathcal{S}_\mathcal{H}}(K)$. If $$\phi(K):=\sup\limits_{f\in\widehat{\mathcal{S}_\mathcal{H}}(K)}\abs{a_2}.$$
Then \begin{enumerate}
\item for $\phi(K)\leq 2K$, we have $$f\in h^p\quad \left(0<p<\frac{1}{2K}\right);$$
\vskip.05in
\item for $\phi(K)>2K$, we have $$f\in h^p\quad \left(0<p<\frac{1}{\phi(K)}\right).$$
\end{enumerate}
Both of these bounds are sharp.}
\end{prop}

\begin{prop} \label{c1}
Let $f\in\widehat{\mathcal{S}_\mathcal{H}}(K,\lambda)$. %such that $f$ preserves affine invariance. % preserve affine and linear invariant. 
Then
\begin{align*}
\sup\limits_{f\in\widehat{\mathcal{S}_\mathcal{H}}(K,\lambda)}\abs{a_2}=\sqrt{1+\frac{\lambda}{2}+\frac{1}{2}\left(\frac{K-1}{K+1}\right)^2}+\frac{K-1}{2\left({K+1}\right)}. 
\end{align*}
\end{prop}

\begin{theorem}\label{t1}
Let $f\in\widehat{\mathcal{S}_\mathcal{H}}(K,\lambda)$ with $K\geq1$ and $\lambda\geq0$. %such that $f$ preserves affine invariance. 
Then
\begin{enumerate}
\item for $0\leq\lambda\leq6$ and $K\geq1$, we have $$f\in h^p\quad \left(0<p<\frac{1}{2K}\right);$$
\vskip.05in
\item for $\lambda>6$ and $K\geq K_1>1$, we have $$f\in h^p\quad \left(0<p<\frac{1}{2K}\right);$$
\vskip.05in
\item for $\lambda>6$ and $1\leq K<K_1$, we have $$f\in h^p\quad \left(0<p<\frac{1}{\varphi(K,\lambda)}\right),$$
\end{enumerate}
where
\begin{equation}\label{1}\varphi(K,\lambda):=\sup\limits_{f\in\widehat{\mathcal{S}_\mathcal{H}}(K,\lambda)}\abs{a_2}
=\sqrt{1+\frac{\lambda}{2}+\frac{1}{2}\left(\frac{K-1}{K+1}\right)^2}+\frac{K-1}{2\left({K+1}\right)},\end{equation}
and 
$K_1$ is the unique solution in $(1,+\infty)$ of the equation
\begin{equation}\label{2}16K^4+24K^3-(2\lambda-11)K^2-2(2\lambda-1)K-2\lambda-5=0.\end{equation}
All of these results are sharp.
\end{theorem}

\begin{proof}[Proof of Theorem \ref{t1}]
If $f\in\widehat{\mathcal{S}_\mathcal{H}}(K,\lambda)$, %preserve affine and linear invariant, 
it follows from Proposition \ref{c1} that
$$\sup\limits_{f\in\widehat{\mathcal{S}_\mathcal{H}}(K,\lambda)}\abs{a_2}=\sqrt{1+\frac{\lambda}{2}+\frac{1}{2}\left(\frac{K-1}{K+1}\right)^2}+\frac{K-1}{2\left({K+1}\right)}\geq1.$$
For convenience, we define a real function $\Phi(K)$ with one parameter $\lambda\, (\lambda\geq0)$ as follows: \begin{equation}\label{33}\Phi(K):=\sqrt{1+\frac{\lambda}{2}+\frac{1}{2}\left(\frac{K-1}{K+1}\right)^2}+\frac{K-1}{2\left({K+1}\right)}-2K\quad(K\geq1).\end{equation}

Now, we divide the proof into three cases.
\vskip.05in
(i) When $0\leq\lambda\leq6$ and $K\geq1$, we know that $\Phi(K)$ is a decreasing continuous function with respect to $K$
in $[1,+\infty)$,
and $$\Phi(1)=\sqrt{1+\frac{\lambda}{2}}-2\leq0,$$ which implies that $$\Phi(K)\leq0\quad(K\geq1).$$ Thus, we have
\begin{equation}\label{3}\varphi(K,\lambda)=\sqrt{1+\frac{\lambda}{2}+\frac{1}{2}\left(\frac{K-1}{K+1}\right)^2}+\frac{K-1}{2\left({K+1}\right)}\leq 2K,\end{equation}
that is, \begin{equation}\label{4}\frac{1}{\varphi(K,\lambda)}\geq\frac{1}{2K}.\end{equation}
By Proposition \ref{prop1},
it shows that $$f\in h^p\quad \left(0<p<\frac{1}{2K}\right).$$
\vskip.05in
(ii) When $\lambda>6$, we know that $$\Phi'(K)<0\quad (K\geq1).$$ Furthermore, we find that $$\Phi(1)=\sqrt{1+\frac{\lambda}{2}}-2>0$$
and $$\Phi\left(\lambda\right)<0.$$ Applying the zero point theorem to the real continuous function $\Phi(K)$ in bounded closed interval
$[1,\, \lambda]$, we deduce that the equation \begin{equation*}\sqrt{1+\frac{\lambda}{2}+\frac{1}{2}\left(\frac{K-1}{K+1}\right)^2}+\frac{K-1}{2\left({K+1}\right)}=2K\quad(K\geq1;\, \lambda>6),\end{equation*} 
which is equivalent to the equation \eqref{2}, exists a unique solution $K_1$ in $(1,+\infty)$.
If $K\geq K_1>1$, we know that
 \eqref{4} also holds. By Proposition \ref{prop1}, it implies that $$f\in h^p\quad \left(0<p<\frac{1}{2K}\right).$$
\vskip.05in
(iii) When $\lambda>6$ and $1\leq K<K_1$, we see that
\begin{equation*}\label{5}\varphi(K,\lambda)=\sqrt{1+\frac{\lambda}{2}+\frac{1}{2}\left(\frac{K-1}{K+1}\right)^2}+\frac{K-1}{2\left({K+1}\right)}>2K,\end{equation*}
or equivalently, $$\frac{1}{\varphi(K,\lambda)}<\frac{1}{2K}.$$ By Proposition \ref{prop1},
it means that $$f\in h^p\quad \left(0<p<\frac{1}{\varphi(K,\lambda)}\right).$$
The proof of Theorem \ref{t1} is thus completed.
\end{proof}

\end{document}
